% Research cards: RED, CON
% Checked on the web on September 23, 2026. "(to be confirmed)" marks what was not verified in a source.

\section{Network operation and continuity}

\subsection*{Shared sources}

\begin{table}[H]
\centering\small
\caption{Sources shared by the RED and CON modules.}
\begin{tabularx}{\textwidth}{P{3.6cm} L P{3.6cm}}
\toprule
\textbf{Source} & \textbf{What it provides} & \textbf{Pitfalls} \\
\midrule
ANEEL, Indicadores Coletivos de Continuidade [1] & measured DEC and FEC and limits by set and month, components, compensation paid, set attributes; since 2000 & limits change at each tariff review; sets are reconfigured \\
ANEEL, Interrupções nas Redes de Distribuição [2] & every interruption since 2017, with \textbf{BDGD feeder and substation code}, set, start and end, cause, reason for exclusion (expurgo), consumers affected; ODbL & cause reported by the distribution utility; does not include permit-holder utilities \\
ANEEL, BDGD [6] and conversion to OpenDSS [7,8] & ANEEL computes technical losses by building the BDGD network in OpenDSS; \emph{bdgd2opendss} (MIT) replicates the methodology; alternatives validated in an undergraduate thesis (4.1\% error in total active power) & static model; \textbf{does not model the consumer's net generation}, which limits reverse power flow studies \\
ONS, constrained-off of wind and photovoltaic plants by plant [10] & half-hourly since 2021, with generation, limited generation, reference, reason for the constraint (REL, CNF, ENE, PAR), origin (local or systemic), connection point; CC-BY & reference generation is an estimate; data are reprocessed \\
ONS, transmission lines and availability [12] & registry of lines above 230 kV; availability of transmission functions by month & registry without history \\
EPE, power flow cases [13] & ANAREDE cases from the PDE and the PAR/PEL & CEPEL format; laborious conversion \\
Lightning [16,17] & RINDAT and BrasilDAT (access by agreement or contract); GLM from GOES, free & GLM does not separate cloud-to-ground from intracloud \\
\bottomrule
\end{tabularx}
\end{table}

\paragraph{PRODIST Module 8 rules (REN 956/2021) used as parameters [4].}
\begin{itemize}
  \item Voltage from 2.3 to 69~kV: adequate between 0.93 and 1.05~pu; precarious
    between 0.90 and 0.93; critical below 0.90 or above 1.05. Low voltage
    220/127~V: adequate from 202 to 231~V and from 117 to 133~V.
  \item DRP and DRC over 1,008 10-minute readings, limits of 3\% and 0.5\%,
    with compensation proportional to the excess.
  \item Compensation for DIC, FIC and DMIC computed with the VRC (TUSD fio B)
    and factors by voltage level; cap of 18 times the VRC; in a simultaneous
    violation only the largest is paid; DICRI separately.
  \item Exclusions from DEC and FEC include interruptions in an emergency
    situation and Critical Day, defined as the day on which the set's
    emergency occurrences exceed the mean plus three standard deviations of
    the previous 24 months [5].
\end{itemize}
Check the version of the module in force before coding the rules [4].

% ---------------------------------------------------------------------------
\subsection{RED \textperiodcentered{} Network operation and planning}

\begin{ficha}{RED-01 \textperiodcentered{} Forecast loading hot spots}{V2}
\campo{Data} Level 0: transformer and substation power rating, monthly
energy by unit and typical curves from the BDGD, converted to OpenDSS [6,7].
Level 1: client billing; level 3: transformer metering.
Admissible overload by the hot-spot thermal model (IEEE C57.91) [42].
\campo{Approaches} Baseline: peak estimated by energy × typical curve
over rated power, projected with load growth. Recommended:
peak quantiles by transformer and feeder, aggregated by the topology,
with output as probability of exceeding the limit and flow confirmed in
OpenDSS. Recent literature: overload alarm classifiers with smart
metering [38].
\campo{Metrics} Precision and recall at top-k against replacements due to
burnout or metering; Brier and calibration of the overload probability; pinball
of the peak.
\campo{Pitfalls} Monthly energy × typical curve smooths the peak; the
overload label only exists where there has already been a problem; outdated
nameplate; load transfers outside the static model.
\end{ficha}

\begin{ficha}{RED-02 \textperiodcentered{} Forecast voltage violation}{V3}
\campo{Data} BDGD electrical model through \emph{bdgd2opendss} [7,8]; forecast
load and MMGD; Module 8 limits [4]; labels generated by time-series
simulation in OpenDSS over scenarios; validation with complaints and sampled
DRP and DRC measurements from the client.
\campo{Approaches} Baseline: full OpenDSS by feeder in the
peak scenario and the maximum-MMGD with minimum-load scenario, feasible in batch.
Recommended: GNN surrogate trained on the simulations of the client's network,
ranking buses and sections, with \textbf{every alert confirmed by OpenDSS}.
Recent literature: PowerFlowNet [26]; PowerFlowMultiNet for unbalanced
three-phase networks, about 100 times faster [27]; physical biases for
voltage in distribution [28]; PF$\Delta$ generalization benchmark [29].
\campo{Metrics} Voltage error in pu; \textbf{recall of violations} above
precision; fraction of alerts confirmed; performance on feeders not
seen in training.
\campo{Pitfalls} Without modeling of net generation, overvoltage from MMGD
needs its own model [7]; regulator and capacitor settings may be
missing; GNN degrades outside the training topologies [29].
\end{ficha}

\begin{ficha}{RED-03 \textperiodcentered{} Prioritization of reinforcement works}{V3}
\campo{Data} Risks from RED-01, CAR-06 and MMG-06; the client's work costs;
taxonomy of the Plano de Desenvolvimento da Distribuição (expansion, improvement,
renewal) [15]; continuity benefits in reais (CON-01, CON-02).
\campo{Approaches} Baseline: greedy ranking by risk avoided per real.
Recommended: multi-year integer programming (knapsack) with dependencies,
mutually exclusive alternatives, annual limit and load and MMGD scenarios, with
the benefit of each work evaluated in OpenDSS. Recent literature: two-stage
stochastic or robust planning [39].
\campo{Metrics} Risk reduced per real; regret in backtest; stability of the
ranking across runs.
\campo{Pitfalls} Benefit depends on uncertain forecasts (MMGD);
regulatory recognition of the investment not modeled (to be confirmed).
\end{ficha}

\begin{ficha}{RED-04 \textperiodcentered{} Constraint and curtailment by region}{V3}
\campo{Data} Constrained-off by plant from the ONS, half-hourly, with reason,
origin and connection point [10], measured by routine RO-AO.BR.13 [11]. For the
map, georeference the connection point through the registry of lines and
substations [12].
\campo{Approaches} Baseline: aggregation of unrealized generation by
connection point, reason and origin, with monthly trend. Recommended: model
of constraint probability and expected energy by region, separating the
energy constraint (systemic) from the electrical ones (local). Recent literature:
reduced network (PTDF) over the EPE cases, calibrated with the history, to
estimate the marginal curtailment of new generation (to be confirmed).
\campo{Pitfalls} Reference generation is an ONS estimate; only the REL reason
goes on to reimbursement at the CCEE; does not cover MMGD.
\end{ficha}

\begin{ficha}{RED-05 \textperiodcentered{} Reconfiguration recommendation}{V4}
\campo{Data} OpenDSS model of the BDGD with switches and normal state; the
client's remote-controlled switches; forecast load; radiality, voltage,
thermal limit and protection constraints.
\campo{Approaches} Baseline: branch exchange of Baran and Wu (1989) [33].
Recommended: mixed-integer conic optimization over DistFlow generates a few
configurations, each one verified in OpenDSS and presented as a suggestion.
Recent literature: reinforcement learning over graphs for restoration in
real time, tested on IEEE networks [34].
\campo{Pitfalls} Protection coordination outside the model; switch state
in the BDGD may not reflect operation; gap between test and real networks.
Shadow mode only.
\end{ficha}

\begin{ficha}{RED-06 \textperiodcentered{} N-1 contingency screening (transmission)}{V4}
\campo{Data} Model of the client's network or ANAREDE cases from the EPE [13];
unavailability of transmission functions [12]; financial impact through the
Parcela Variável por Indisponibilidade (REN 729/2016) [14].
\campo{Approaches} Baseline: distribution factors (PTDF, LODF) and full
AC on the most severe ones. Recommended: GNN calibrated for recall close to
100\%, with full AC on everything above a low threshold. Recent literature: PINN
with GNN for N-1 to N-3, about 400 times faster, tested up to 118 buses
[31]; input-convex networks with a reliability guarantee [32].
\campo{Pitfalls} Since a false negative is unacceptable, the model only orders
the cases for the full AC. The
literature goes up to 118 buses, much smaller than the SIN; the official
security assessment belongs to the ONS.
\end{ficha}

% ---------------------------------------------------------------------------
\subsection{CON \textperiodcentered{} Continuity and quality}

\begin{ficha}{CON-01 \textperiodcentered{} Forecast DEC and FEC by set}{V2}
\campo{Data} DEC and FEC with limits [1]; interruptions by event linked to the
feeder since 2017 [2]; emergency occurrences [3]; weather (INMET,
reanalysis) and lightning.
\campo{Approaches} Baseline: seasonal naive and ETS by set.
Recommended: monthly global model by set over the regulated components,
with the exclusions modeled separately; output as quantiles and probability of
exceeding the annual limit. Recent literature: frequency × duration by feeder
in Monte Carlo over weather scenarios, reconciled up to the set (to be
confirmed). In Brazil: neural networks for indicators [41] and ML in the
definition of set limits [40].
\campo{Metrics} MAE and pinball of DEC and FEC; Brier of the violation
probability; hit rate of the ranking of sets at risk.
\campo{Pitfalls} The Critical Day uses a 24-month moving threshold and makes
the exclusion dependent on the weather [5]; reconfiguration of sets breaks
series.
\end{ficha}

\begin{ficha}{CON-02 \textperiodcentered{} DIC and FIC violation risk (compensation)}{V2}
\campo{Data} Module 8 formulas and factors [4]; compensation aggregated by
set in the open data [1]. \textbf{Monthly DIC and FIC by unit are
in the BDGD}, in the UCBT and UCMT tables, verified in the files: sizing the
exposure is \textbf{level 0}. Caveat: the BDGD has a lag of 8 to 20 months and
the unit is anonymized by hash, so the public data serve to size
and validate the method, not to track the current quarter --- that is
level 1, with the distribution utility's code mapping.
\campo{Approaches} Baseline: historical violation rate and mean
compensation by feeder. Recommended: simulate DIC, FIC and DMIC by unit
from an event model by section, apply the compensation rules
(maximum, cap) and aggregate the expected value by feeder or transformer.
\campo{Metrics} Compensation error in reais by feeder and quarter;
precision at top-k of the ranking.
\campo{Pitfalls} Monthly, quarterly and annual limits at the same time;
nonlinear rules; LGPD on the data by unit.
\end{ficha}

\begin{ficha}{CON-03 \textperiodcentered{} Interruption risk by weather event}{V3}
\campo{Data} Labels: interruptions with cause and exclusion, and emergency
events [2]. Weather forecast (CLI-01), stroke density from
CLI-03 [16,17], vulnerability (vegetation from ATV-03, network age and type from
the BDGD).
\campo{Approaches} Baseline: count regression (Poisson or negative
binomial) by cell and day over gust, rain and lightning. Recommended: the
pattern of the UConn model, in operation since 2015: tree ensemble by
event on a 2~km grid, with weather, vegetation and
infrastructure variables and modules by weather type [19--21]. Recent literature:
storms as tracked objects [20]; two stages with open data
[24]; GNN with tree risk from LiDAR [25].
\campo{Metrics} Error of the total interruptions by event; spatial POD, FAR
and CSI; skill by lead time. \textbf{Leave-one-event-out and
leave-one-region-out validation.}
\campo{Pitfalls} Random splits inflate the performance estimate: with spatial
and temporal validation, many models do not beat the baseline [23]. Weather
forecast error propagates. Transfer of models from temperate climates
to tropical convection is not proven.
\end{ficha}

\begin{ficha}{CON-04 \textperiodcentered{} Probable fault location}{V4}
\campo{Data} Recloser and relay events (sequence of events, current,
oscillography), last signal from the smart meters, customer calls,
OpenDSS topology. Requires level 3 for real-time use.
\campo{Approaches} Baseline: topological inference (common ancestor of the
units without power, downstream of the device that operated) combined with the
impedance method in OpenDSS. Recommended: physical candidates ranked
by a classifier with cause priors (vegetation, lightning, history).
Recent literature: spatio-temporal GNN on partially observed networks [36];
digital twin with low-voltage meters [37].
\campo{Metrics} Section hit rate at top-1 and top-3; error in km; reduction in
kilometers patrolled.
\campo{Pitfalls} Scarce real labels; multiple faults in a storm; DG
distorts the impedance method.
\end{ficha}

\subsection*{CON module diagrams}
\addcontentsline{toc}{subsection}{CON module diagrams}

The module has two rhythms: the \textbf{regulatory} one, monthly, over sets and
feeders (CON-01, CON-02), and the \textbf{event} one, from hours to days,
triggered by each weather forecast run (CON-03). The diagrams show
on which entities each layer is drawn, who uses what, the states of the two
central objects (the weather event and the set in the regulatory year) and the
flow that links data, models and decisions.

\begin{figure}[H]
\centering
\includegraphics[width=\textwidth]{con-camadas-en.png}
\caption{CON: layers by entity and the data that feed them. Level 0, except
CON-04; CON-02 requires level 1 for the current quarter.}
\end{figure}

\begin{figure}[H]
\centering
\includegraphics[width=0.78\textwidth]{con-casos-de-uso-en.png}
\caption{CON: use cases, separated by the module's two rhythms.}
\end{figure}

\begin{figure}[H]
\centering
\begin{minipage}[t]{0.40\textwidth}\centering
\includegraphics[width=\linewidth]{con-estados-evento-en.png}
\end{minipage}\hfill
\begin{minipage}[t]{0.56\textwidth}\centering
\includegraphics[width=\linewidth]{con-estados-conjunto-en.png}
\end{minipage}
\caption{CON: state machines of the weather event (left) and of the electrical
set in the regulatory year (right). Labeled events, including false
alarms, return as training labels.}
\end{figure}

\begin{figure}[H]
\centering
\includegraphics[width=\textwidth]{con-fluxo-en.png}
\caption{CON: flow of the monthly cycle, the event cycle and the closing, with
leave-one-event-out and leave-one-region-out validation.}
\end{figure}

\subsection*{References for this section}
{\small
\begin{enumerate}[label={[\arabic*]},leftmargin=2.2em,itemsep=1pt]
\item ANEEL, Indicadores Coletivos de Continuidade: \url{https://dadosabertos.aneel.gov.br/dataset/indicadores-coletivos-de-continuidade-dec-e-fec}
\item ANEEL, Interrupções nas Redes de Distribuição: \url{https://dadosabertos.aneel.gov.br/dataset/interrupcoes-de-energia-eletrica-nas-redes-de-distribuicao}
\item ANEEL, Ocorrências Emergenciais: \url{https://dadosabertos.aneel.gov.br/dataset/ocorrencias-emergenciais-nas-redes-de-distribuicao}
\item ANEEL, PRODIST Module 8 (REN 956/2021): \url{https://www.gov.br/aneel/pt-br/centrais-de-conteudos/procedimentos-regulatorios/prodist}
\item REN ANEEL 956/2021, PRODIST Module 1 (Critical Day)
\item ANEEL, BDGD: \url{https://dadosabertos.aneel.gov.br/dataset/base-de-dados-geografica-da-distribuidora-bdgd}
\item Radatz, bdgd2opendss: \url{https://github.com/PauloRadatz/bdgd2opendss}
\item Aderaldo (2023), BDGD-Tools, UFC: \url{https://repositorio.ufc.br/handle/riufc/75567}
\item[] \refstepcounter{enumi}
\item ONS, constrained-off by plant: \url{https://dados.ons.org.br/dataset/restricao_coff_eolica_detail}
\item ONS, RO-AO.BR.13, constrained-off measurement
\item ONS Dados Abertos, lines and availability: \url{https://dados.ons.org.br/dataset/linha-transmissao}
\item EPE, transmission databases: \url{https://www.epe.gov.br/pt/leiloes-de-energia/leiloes-de-transmissao/bases-de-dados}
\item ANEEL, REN 729/2016 (Parcela Variável)
\item ANEEL, Plano de Desenvolvimento da Distribuição: \url{https://www.gov.br/aneel/pt-br/assuntos/distribuicao/plano-de-desenvolvimento-da-distribuicao}
\item INPE/ELAT, RINDAT and BrasilDAT: \url{http://www.inpe.br/webelat/rindat/}
\item NASA Earthdata, GLM: \url{https://www.earthdata.nasa.gov/data/instruments/glm}
\item INMET, BDMEP: \url{https://bdmep.inmet.gov.br/}
\item Wanik et al.\ (2015), \emph{Natural Hazards} 79, doi:10.1007/s11069-015-1908-2
\item Cerrai et al.\ (2019), \emph{IEEE Access} 7
\item UConn Outage Prediction Model, AGU (2021)
\item Extratropical storms outages, \emph{NHESS} 21 (2021)
\item Essus, Vatsavai and Rachunok (2026), leakage in interruption models: \url{https://arxiv.org/abs/2608.24665}
\item Stanishevska and Guikema (2025): \url{https://arxiv.org/abs/2510.03959}
\item GNN and tree risk from LiDAR, \emph{RESS} (2025)
\item Lin et al.\ (2023), PowerFlowNet: \url{https://arxiv.org/abs/2311.03415}
\item Ghamizi et al.\ (2024), PowerFlowMultiNet: \url{https://arxiv.org/abs/2403.00892}
\item Okoyomon, Yaniv and Goebel (2025): \url{https://arxiv.org/abs/2509.25158}
\item PF$\Delta$ (2025): \url{https://arxiv.org/abs/2510.22048}
\item SafePowerGraph (2024): \url{https://arxiv.org/abs/2407.12421}
\item Nakiganda and Chatzivasileiadis, GNN for contingencies: \url{https://arxiv.org/abs/2310.04213}
\item N-k with input-convex networks (2024): \url{https://arxiv.org/abs/2410.00796}
\item Baran and Wu (1989), \emph{IEEE Trans.\ Power Delivery} 4(2)
\item Jacob et al.\ (2024), \emph{Nature Communications}
\item PI-GNN for reconfiguration (2023): \url{https://arxiv.org/abs/2310.00728}
\item Karabulut, Manna and Develder (2026): \url{https://arxiv.org/abs/2604.20403}
\item Fault location in LV with a digital twin, \emph{Energies} 16 (2023), doi:10.3390/en16237850
\item Overload alarm in distribution transformers (2024)
\item Investment in distribution networks with correlation between projects, \emph{Frontiers in Energy Research} (2021)
\item Interruption limits in Brazil with AI, \emph{Energies} 16 (2023), doi:10.3390/en16197012
\item Prediction of quality indicators with neural networks, SBSE
\item IEEE C57.91, loading guide
\end{enumerate}}
